\section{问题与改进}

\subsection{已完成的功能}

\begin{enumerate}
  \item \textbf{保存中间态。} 项目在每次Oracle和Diffusion后保存statevector，形成完整的步骤序列。
  \item \textbf{同时计算概率幅和概率。} 概率幅图给出Oracle后的符号翻转，概率图给出Diffusion后的测量概率提升。
  \item \textbf{生成3D步进图。} 3D图把“基态--步骤--概率”放在同一个坐标系中，记录目标态概率柱随步骤升高的过程。
  \item \textbf{提供交互参数。} 用户可以切换量子比特数、目标态、迭代次数，并通过滑块选择当前步骤。
  \item \textbf{保留量子电路信息。} 除图表外，界面还显示Qiskit标准电路图和交互式操作表，使每一步状态变化对应到具体量子门。
\end{enumerate}

\subsection{不足与改进方向}

本项目按照课程项目要求支持2到3个量子比特。若扩展到更多量子比特，柱状图数量会迅速增加，3D图也会变得拥挤，需要改用筛选、聚类或只突出目标态和若干代表性非目标态的方式。

另一方面，当前噪声模型只是简化的depolarizing noise和readout noise，不代表真实量子硬件的校准噪声。后续可以接入真实后端参数，或加入多目标Grover搜索。
